\section{Introduction}
\label{sec:introduction}

\IEEEPARstart{V}{isual} speech recognition converts silent mouth motion into text and enables communication when audio is unavailable, private, or unreliable~\cite{potamianos2003recent,ma2022visual}. Modern systems combine a spatiotemporal visual front end with a sequence encoder and CTC~\cite{graves2006ctc} or attention-based decoding. Their reported accuracy, however, normally assumes a fixed model. A deployed recognizer instead observes a stream: speakers recur, camera conditions drift, and a user may occasionally correct a transcript. The practical question is therefore not only whether a source model generalizes, but whether a small amount of feedback can improve future predictions without replaying source data or corrupting previously useful behavior.

Speaker adaptation has clear precedent in VSR. Prompt- and LoRA-based personalization can adapt both visual and language components to a target user~\cite{yeo2024personalized}. Continual test-time adaptation (CTTA) similarly studies non-stationary unlabeled streams, including parameter-efficient adaptation, recovery mechanisms, and expert routing~\cite{niu2022eata,wang2022cotta,lee2024becotta,cui2025paint}. These settings do not directly answer our question. VSR personalization is commonly organized as an offline support-set stage, while generic CTTA is primarily evaluated on image classification or semantic segmentation. Sparse human corrections introduce a different resource: labels are direct but intermittent and may themselves be noisy, and every result must be recorded before its label is used.

This paper studies that setting with a frozen Mandarin VSR backbone and lightweight residual adapters. The strict causal path uses one adapter and separates prediction from any later correction. We also audit an earlier dynamic-expert prototype under matched streams, but its known feedback-location flag could affect pre-prediction growth; all expert numbers are therefore retained only as pre-fix diagnostics. This distinction matters because a more complex structure is useful only if it routes domains meaningfully and improves either accuracy or revisit forgetting beyond a single adapter.

Our contributions are empirical and methodological:
\begin{enumerate}
  \item We define an auditable sparse-feedback protocol for continual Mandarin VSR. Predictions are scored before adaptation, stream order and feedback positions are fixed, manifests and checkpoints are hash-locked, and comparisons use utterance-paired bootstrap intervals.
  \item We separate two measured gains: no-feedback online adapter adaptation improves seed-42 CER by 1.72 points, while one correction every ten utterances produces a stable three-seed total gain of 2.42 points over the static model.
  \item We provide a negative diagnostic for the current expert implementation under two routing regimes: the original expert bank collapses almost entirely to one expert, while a validation-calibrated repair over-fragments into eight experts and performs materially worse than one adapter. These pre-fix runs do not establish that dynamic experts are intrinsically ineffective.
  \item We audit where correction gradients should act. Greedy correct-span preservation is worse than whole-sequence replay. Target-conditioned CTC localization significantly beats same-mass randomized support, yet remains significantly worse than replay and is rejected by its validation gate.
\end{enumerate}

The paper deliberately separates supported claims from open work. The validation-only router gate is resolved as a no-go, so no repaired-router test run or additional expert seed is reported. An adapter-only entropy baseline is included as a failure diagnostic; stronger matched CTTA baselines remain required before a method-comparison submission.
